% SYNTHETIC TEST FIXTURE for paper-copilot, generated by scripts/gen-large-fixture.mjs (seed 20261106).
% A real-size, multi-file revtex4-2 paper: every number of the prose is in ../repo/results or
% ../repo/docs/definitions.md. Regenerate with `node scripts/gen-large-fixture.mjs`; do not edit by hand.
\documentclass[aps,prd,preprint,superscriptaddress,nofootinbib]{revtex4-2}
\usepackage{amsmath}
\usepackage{graphicx}
\usepackage{hyperref}
\newcommand{\ttbar}{\ensuremath{t\bar{t}}}

\begin{document}

\title{Fixture: displaced tracks from a graph tracker trained without signal}
\author{A. Fixture}
\affiliation{Nowhere Institute}
\author{B. Coauthor}
\affiliation{Elsewhere Laboratory}

\begin{abstract}
Charged particles from the decay of a long-lived particle start far from the collision point.
Trackers tuned on prompt particles lose many of them, and the classical remedy, a second tracking pass, is slow.
We train a graph tracker on generic top-quark events only and test it on three signal samples that it never saw.
On displaced targets, it finds 90.06\% of the targets, compared with 76.52\% for a classical baseline, at a fake-track rate of 0.75\% against 1.63\%.
The gain grows with the production radius and does not depend on the signal model.
\end{abstract}

\maketitle

% Introduction (one sentence per line).
\section{Introduction}
\label{sec:intro}

Many extensions of the Standard Model predict particles that live long enough to travel a measurable distance before they decay~\cite{Dubois2008LLP,Vargas2024HNL}.
Their charged decay products start far from the beam line, so their tracks do not point back to the collision point~\cite{Kowalski2013Displaced}.
Searches for such signatures have become a central part of the physics programme of the LHC experiments.
Heavy neutral leptons and dark photons are two well-motivated examples~\cite{Vargas2024HNL,Chen2012DarkPhoton}.
In both cases the signal is a few tracks that appear in the middle of the detector.
Finding these tracks is the first step of every such search, and every later step inherits its losses.

Standard track reconstruction starts from seeds near the beam line and extends them outward with a Kalman filter~\cite{Hansen2014Kalman,Okafor2019Combinatorial}.
This strategy is efficient for prompt particles and fast enough for the trigger.
It loses particles whose first hit lies beyond the innermost layers, because no seed can be formed for them.
Dedicated large-radius tracking recovers part of these particles at a large cost in time and fake tracks~\cite{Rossi2018LRT}\footnote{Large-radius passes usually run only on the hits that the prompt pass left unused.}.
The cost grows with pile-up, which will increase several-fold at the high-luminosity LHC~\cite{Jensen2017Pileup,Okafor2013Upgrade}.
A reconstruction that treats prompt and displaced particles alike would remove this trade-off.

Early searches for long-lived particles relied on hand-written selections of displaced vertices.
Their reach was limited by the tracking, because a vertex can only be built from tracks that were found.
Later searches added dedicated passes of track finding with relaxed requirements on the impact parameter.
These passes raised the efficiency for displaced particles but also the number of fake tracks, which then had to be removed by tighter vertex selections.
The net gain depended strongly on the signal model, because the relaxed passes were tuned for particular lifetimes.

Graph neural networks offer such a reconstruction~\cite{Moreau2016GNNReview}.
They represent the hits of an event as the nodes of a graph and the possible track segments as its edges.
A network then learns which edges belong to the same particle~\cite{Fischer2011Interaction,Lindqvist2020MessagePassing}.
The Exa.TrkX pipeline showed that this approach reaches high efficiency on the TrackML detector~\cite{Yilmaz2023ExaTrkX,Tanaka2021TrackML}.
Because the graph does not assume that a particle comes from the beam line, displaced particles are not excluded by construction.
Recent studies extended these pipelines to large-radius tracks and found efficiencies close to those of prompt tracks~\cite{Sato2008EdgeClass,Zhang2013Metric}.

Learned tracking has developed quickly over the last decade.
The first studies classified hit pairs with simple networks and built tracks by hand.
Graph networks then made it possible to reason over many hits at once.
Metric learning replaced hand-written geometric rules in the construction of the graph.
End-to-end pipelines now take raw hits and return tracks with parameters.
Throughout this development, the training events were drawn from the process under study.

Learned trackers change this picture in two ways.
They can be trained to find displaced and prompt particles in one pass, without separate settings.
They can also be retrained cheaply when the detector or the run conditions change.
Both properties are attractive for a search programme that covers many models.
Both, however, rest on the assumption that the training events teach the network something general about tracks.

However, these studies trained the network on the signal they later evaluate.
A tracker trained on one model of new physics may favour that model and miss others.
In a search this is a bias: the selection would depend on the hypothesis being tested.
The machine-learning literature calls the question out-of-distribution generalization~\cite{Tanaka2015Generalization,Abe2020Domain}.
For tracking it has not been answered, because the training and test samples were always drawn from the same process.
We ask whether a tracker trained only on Standard Model events can find displaced particles as well as one trained on the signal.

This paper tests that assumption directly.
We do not propose a new architecture, and we do not tune the pipeline for displaced particles.
We take an established pipeline, keep every part of it fixed, and change only what it is trained on.
The comparison is therefore a measurement of what the training sample contributes.
A positive answer would mean that experiments can train one tracker on the abundant generic events they already simulate.

To answer it, we train two copies of the same pipeline.
One copy sees only top-quark pair events, which contain almost no displaced particles.
The other copy sees the displaced signal it will be tested on.
Both copies share the graph construction, the network, the track builder and every hyperparameter.
We then compare them on displaced particles generated after the design of the study was frozen.
A tracker that wins this comparison does not need the signal to be trained.

The answer is yes: the generic-trained tracker finds 2.49 points more displaced targets than the signal-trained one.
It also produces fewer fake tracks and fewer duplicates~\cite{Jensen2011Calibration}.
The result holds across particle species, displacement and pile-up.
It does not depend on the choice of training seed or on the bootstrap of the test events~\cite{Novak2023Bootstrap,Usman2011Seeds}.
The rest of this paper explains how we reached it.

Our contribution is a controlled measurement, not a new method.
It isolates the effect of the training sample on a learned tracker.
It reports that effect with an uncertainty that includes the randomness of the training.
And it checks that the effect survives changes of the signal model, the conditions and the analysis choices.

Section~\ref{sec:method} describes the detector, the samples, the graph and the network.
Section~\ref{sec:results} compares the two trackers and shows where the difference comes from.
Section~\ref{sec:summary} summarizes the result, and Appendix~\ref{app:stat} gives the details of the uncertainty.
The figures that carry the evidence are Fig.~\ref{fig:eff_disp} for the displacement and Fig.~\ref{fig:pileup} for pile-up.

Our study builds on a long line of work on fast tracking~\cite{Xu2010Hough,Esposito2015Cellular,Quinn2022Trigger}.
Graph-based pipelines have been ported to accelerators and now run within the time budget of a software trigger~\cite{Ito2021GPU,Petrov2009Inference}.
Their robustness to detector conditions has been studied for prompt particles~\cite{Hansen2008Robustness}.
Displaced particles have received less attention, because they are rare in the samples used for training.
This scarcity is exactly what makes generic training attractive: the network cannot learn the signal it never sees\footnote{Heavy-flavour hadrons travel a few millimetres before they decay, so the generic sample does contain some displaced tracks.}.

Filtering the candidate edges before classification is a common way to reduce the size of the graph~\cite{Garcia2018Filtering}.
We do not filter, because a filter trained on one sample could itself introduce the bias we want to measure.
The graph is therefore larger than in pipelines optimized for speed, and the edge classifier sees every candidate edge.
This choice makes the comparison cleaner at the cost of some time per event.

We simulate collisions at a centre-of-mass energy of 14~TeV with a generic all-silicon detector~\cite{Abe2009ACTS,Dubois2019Pythia}.
The detector response is simulated with a fast model that keeps the material and the magnetic field of the full geometry~\cite{Weber2014Geant,Kowalski2024Delphes}.
The study therefore measures the tracker, not the detector, and its numbers hold within this simulation only~\cite{Yilmaz2017Material,Rossi2012Resolution}.
We expect the comparison between the two trackers to carry over to a real detector, because both suffer the same effects.

The question matters beyond the LHC, where future colliders will face higher rates and new long-lived signatures~\cite{Chen2023FCC,Vargas2018Software}.
It also matters for model-independent searches, which need a reconstruction that does not encode a signal hypothesis~\cite{Moreau2010Anomaly}.
Statistical tools for such searches assume that the selection does not depend on the model being tested~\cite{Quinn2016Statistics}.
A generic-trained tracker satisfies this assumption by construction.

We keep the comparison fair in three ways.
First, both trackers use the same architecture and the same number of training events.
Second, every choice of the pipeline was fixed on validation events before the test events were examined.
Third, the test particles come from a signal sample that neither tracker was trained on.
With these precautions, a difference between the trackers can only come from what they were trained on.

% Method (one paragraph per line).
\section{Detector, samples and method}
\label{sec:method}

\subsection{Detector and samples}
\label{sec:detector}

The detector has 22 silicon layers arranged in a barrel and two endcaps, inside a solenoidal magnetic field. The innermost layers are pixel sensors, and the outer layers are strip sensors with a coarser resolution along the beam axis. Figure~\ref{fig:detector} shows the layout in the transverse and longitudinal views. Each charged particle leaves one hit per layer it crosses, plus occasional extra hits where modules overlap. Hits are recorded as three-dimensional space points in cylindrical coordinates.

\begin{figure}[tbp]
  \centering
  \fbox{\parbox[c][0.45\linewidth][c]{0.8\linewidth}{\centering placeholder artwork: detector}}
  \caption{Layout of the simulated detector in the transverse (left) and longitudinal (right) views. Pixel layers are drawn in dark grey and strip layers in light grey.}
  \label{fig:detector}
\end{figure}

We use three samples of simulated events. The generic sample contains 10,000 top-quark pair events without pile-up. The signal sample contains 10,000 events in which a heavy neutral lepton decays to a muon and two other charged particles\footnote{The decay length of each heavy neutral lepton is drawn from an exponential distribution with the nominal lifetime.}. The external test sample was generated with a different lepton mass and lifetime after the design of the study was fixed. Each sample is split into 8,000 training, 1,000 validation and 1,000 test events; Table~\ref{tab:samples} lists them. Selection of anything (models, thresholds, working points, bins) uses validation events only.

Events are generated with a standard parton-shower generator and passed through a fast simulation of the detector. The simulation propagates each charged particle through the magnetic field and the material of every layer. Multiple scattering and energy loss are applied layer by layer, and hits are smeared with the resolution of the sensor type. Noise hits are not simulated, because their rate in silicon detectors is small compared with the hits from particles. Secondary particles from interactions with the material are kept and treated like any other particle. They are a realistic source of displaced tracks in the generic sample.

The heavy neutral lepton is produced in the decay of a W boson together with a prompt muon. It travels a distance set by its lifetime and then decays to a second muon and a pair of quarks or leptons. The charged decay products therefore start at the position of the decay, which can be anywhere inside the tracker. We use the prompt muon of the same event as a reference that both trackers should find. Its efficiency is close to one for both trackers, which shows that the comparison is not affected by a gross failure of either.

Before the graph is built, the hits of each event are sorted by layer and by azimuthal angle. Hits from the same module that lie closer than the sensor pitch are merged into one space point. The merged space point keeps the identity of the particle that left most of its charge. Space points are written to disk together with the truth information needed for training and evaluation. The truth information never reaches the network at inference time.

Events are processed one at a time, without batching across events at inference. Each event is independent, so the pipeline can be distributed over many processors without communication. The same code runs the training, the validation and the evaluation, with only the input files changed. This guarantees that the two trackers see the same preprocessing.

\begin{table}[tbp]
  \caption{Simulated samples and their splits. The external test sample is used only for the final comparison.}
  \label{tab:samples}
  \centering
  \begin{tabular}{lrrr}
    \hline\hline
    Sample & Training & Validation & Test \\
    \hline
    Top-quark pairs & 8,000 & 1,000 & 1,000 \\
    Heavy neutral lepton & 8,000 & 1,000 & 1,000 \\
    External signal & --- & --- & 1,000 \\
    \hline\hline
  \end{tabular}
\end{table}

A target particle has $p_T > 1$~GeV, $|\eta| < 4$ and at least five space points. A displaced particle is a target with a transverse impact parameter $|d_0| > 2$~mm. The impact parameter is defined from the particle's own hits by a circle fit in the transverse plane,
\begin{equation}
  d_0 = \left(\vec{p} \times \hat{t}\right)_z ,
  \label{eq:d0}
\end{equation}
where $\vec{p}$ is the point of closest approach to the beam axis and $\hat{t}$ the unit direction of motion at that point.

\subsection{Graph and network}
\label{sec:graph}

The pipeline has four stages: graph construction, edge classification, edge selection and track building. Figure~\ref{fig:pipeline} shows the stages and the information that flows between them. Graph construction connects each hit to the hits on adjacent layers that a particle could plausibly reach. A module map, learned from the training events, lists the pairs of modules that particles actually connect\footnote{The module map is built from training events only, so it contains no information about the test sample.}. Only hit pairs with $|\Delta z| < 300$~mm on adjacent layers enter the graph. On average a graph holds 43,260 nodes and 146,892 edges per event. It keeps 98.23\% of the true edges at an edge purity of 0.42.

The edge classifier is an interaction network that updates node and edge features in alternating steps. Each of the 8 message-passing steps lets information travel one edge further along a candidate track. The node and edge networks are multilayer perceptrons with 128 hidden units each. The full model has 1,344,392 trainable parameters. After the last step, a final perceptron assigns each edge a score between zero and one.

\begin{figure*}[tbp]
  \centering
  \fbox{\parbox[c][0.45\linewidth][c]{0.9\textwidth}{\centering placeholder artwork: pipeline}}
  \caption{The four stages of the pipeline. Graph construction and track building are the same for both trackers; only the training events of the edge classifier differ.}
  \label{fig:pipeline}
\end{figure*}

The network is trained with a weighted binary cross-entropy,
\begin{equation}
  \mathcal{L} = -\sum_{e} w_e \left[ y_e \log s_e + (1 - y_e) \log (1 - s_e) \right] ,
  \label{eq:loss}
\end{equation}
where $s_e$ is the score of edge $e$, $y_e$ is one for a true edge and zero otherwise, and the weight $w_e$ balances the rare true edges against the many false ones.

\subsection{Training, track building and metrics}
\label{sec:metrics}

We train for 60 epochs with the Adam optimizer~\cite{Bauer2016Adam}, a batch of 16 events and a learning rate of $3\times10^{-4}$. Every configuration is trained with 3 random seeds, and the reported numbers are averages over the seeds. The two trackers differ only in their training events: generic events for one, signal events for the other. Training stops at the epoch with the best validation efficiency, judged on validation events of the training process itself. No test event influences any choice of the pipeline.

Edges with a score above one half are kept. The threshold was fixed on validation events of the generic sample and is the same for both trackers. Track building then follows the kept edges from each hit with no incoming edge to the hit with no outgoing edge. Where a path branches, the builder follows the edge with the higher score and starts a new candidate at the other. A candidate becomes a track when it holds at least five hits.

A reconstructed track matches a particle if more than half of its hits come from that particle and more than half of the particle's hits are on the track. This criterion is called double majority. The efficiency is the fraction of target particles matched by a track,
\begin{equation}
  \epsilon = \frac{N_{\text{matched}}}{N_{\text{targets}}} ,
  \label{eq:eff}
\end{equation}
where both counts use the same population of targets; Eq.~(\ref{eq:eff}) is evaluated separately for each tracker.

A fake track is a track with at least five hits that matches no particle\footnote{A track may match at most one particle and a particle at most one track under double majority.}. The fake-track rate is the fraction of such tracks among all tracks with at least five hits. A duplicate is a second track matched to a particle that already has one. The three numbers together describe what a search would see: the particles it finds, the noise it must reject and the double counting it must undo.

The headline number is the efficiency difference between the trackers on displaced targets,
\begin{align}
  \Delta &= \epsilon_{\text{generic}} - \epsilon_{\text{signal}} , \label{eq:delta} \\
  \sigma_\Delta^2 &= \sigma_{\text{boot}}^2 + \sigma_{\text{seeds}}^2 . \label{eq:sigma}
\end{align}
with the two components defined in Appendix~\ref{app:stat}.

Its uncertainty combines two sources in quadrature: the bootstrap of the test events and the spread over training seeds~\cite{Novak2023Bootstrap}. The bootstrap draws the same resampled event lists for both trackers, so their difference is not affected by which events happen to be in the test set~\cite{Quinn2016Statistics}. Appendix~\ref{app:stat} gives the details of both components. Both trackers are evaluated on exactly the same test events and the same target particles.

\subsection{Track fit and validation}
\label{sec:validation}

After track building, every track is fitted with a helix to obtain its parameters. The fit uses only the hits on the track and does not depend on which tracker built it. A track whose fit fails is kept for the efficiency but marked in the resolution study. Fits fail mostly for tracks that combine hits of two particles, which are also the tracks that fail double majority. The fitted parameters enter only the resolution study in the results.

We validated every stage of the pipeline on validation events before looking at the test events. The graph construction was checked by counting the true edges it keeps, layer pair by layer pair. The edge classifier was checked by comparing its score distribution for true and false edges between training and validation events. The track builder was checked on graphs that contain only true edges, where it must reproduce the true particles. All checks passed with the same settings for both trackers. We then froze the pipeline, recorded its configuration and only afterwards evaluated the test events.

Training and inference run on one graphics processor per job. A full training run of either tracker takes less than a day. The evaluation of all test samples for one network takes a few hours. All runs use the same software versions, which are recorded with their outputs. Floating-point results differ slightly between processors, but by far less than the seed spread.

% Results (hard-wrapped at about 100 columns).
\section{Results}
\label{sec:results}

On the displaced test particles, the generic-trained tracker finds 95.52\% of the targets. The
signal-trained tracker finds 93.03\% of the same targets. The difference is 2.49~$\pm$~0.24 points,
in favour of the tracker that never saw the signal\footnote{The quadrature sum assumes that the two
components are independent, which holds because they come from different sources of randomness.}.
The test set holds 13,563 target particles, so the difference is far more than a handful of
particles. Table~\ref{tab:headline} lists the efficiencies together with the fake-track and
duplicate rates.

\begin{table}[tbp]
  \caption{Efficiency, fake-track rate and duplicate rate of the two trackers on displaced targets.}
  \label{tab:headline}
  \centering
  \begin{tabular}{lcc}
    \hline\hline
     & Generic-trained & Signal-trained \\
    \hline
    Efficiency (\%) & 95.52 & 93.03 \\
    Fake-track rate (\%) & 0.48 & 1.32 \\
    Duplicate rate (\%) & 1.3 & 2.6 \\
    \hline\hline
  \end{tabular}
\end{table}

The generic-trained tracker also produces fewer fake tracks: 0.48\% against 1.32\%. Its duplicate
rate is 1.3\%, compared with 2.6\% for the signal-trained tracker. A tracker that finds more
particles usually pays with more fake tracks, so the lead on both metrics is not a trade-off.
Figure~\ref{fig:fake} shows the fake-track rate as a function of the track length. The difference in
fake tracks comes mostly from short tracks of five or six hits.

\begin{figure}[tbp]
  \centering
  \fbox{\parbox[c][0.45\linewidth][c]{0.8\linewidth}{\centering placeholder artwork: fake}}
  \caption{Fake-track rate versus the number of hits on the track. Most fake tracks of the signal-trained tracker are short.}
  \label{fig:fake}
\end{figure}

\subsection{Displacement, species and pile-up}
\label{sec:dependence}

Figure~\ref{fig:eff_disp} shows the efficiency as a function of the transverse displacement of the
production vertex. Close to the beam line, between zero and 2~mm, the two trackers find 98.87\% and
98.28\% of the targets. In the bin from 2 to 5~mm, the efficiencies are 98.37\% and 97.16\%. For
production radii from 5 to 20~mm, the efficiencies are 97.65\% for the generic-trained tracker and
95.89\% for the baseline. Between 20 and 50~mm, the generic-trained tracker finds 97.10\% of the
targets and the signal-trained one 94.78\%. The gap between the trackers widens with displacement.

\begin{figure*}[tbp]
  \centering
  \fbox{\parbox[c][0.45\linewidth][c]{0.9\textwidth}{\centering placeholder artwork: eff disp}}
  \caption{Efficiency versus the transverse displacement of the production vertex for displaced targets. Solid: the generic-trained tracker; dashed: the signal-trained tracker. The lower panel shows their difference in points.}
  \label{fig:eff_disp}
\end{figure*}

In the bin from 50 to 100~mm, the efficiencies are 95.62\% and 92.62\%. For production radii from
100 to 200~mm, the efficiencies are 94.26\% for the generic-trained tracker and 90.50\% for the
baseline. Between 200 and 300~mm, the generic-trained tracker finds 91.76\% of the targets and the
signal-trained one 87.28\%. In the bin from 300 to 600~mm, the efficiencies are 89.57\% and 83.93\%.
Beyond the outermost pixel layer, both trackers lose particles that leave fewer than five hits. This
loss is a property of the detector, not of the trackers, and it is the same for both. The widening
gap shows that the generic-trained tracker generalizes better precisely where the signal-trained one
should have its advantage.

Figure~\ref{fig:eff_species} breaks the efficiency down by particle species. For muons, the
efficiencies are 97.52\% and 95.61\%, from 2,289 targets. The 1,754 electrons are found with
efficiencies of 97.52\% and 96.16\%. For charged pions, the efficiencies are 96.23\% and 94.76\%,
from 4,843 targets. The 1,533 charged kaons are found with efficiencies of 96.36\% and 94.06\%. For
protons, the efficiencies are 94.04\% and 92.00\%, from 889 targets. Table~\ref{tab:species} lists
the numbers with their bootstrap uncertainties. The lead holds for every species. Electrons are the
hardest case for both trackers, because they lose energy in the material and their tracks bend more
than the network expects.

\begin{figure}[tbp]
  \centering
  \fbox{\parbox[c][0.45\linewidth][c]{0.8\linewidth}{\centering placeholder artwork: eff species}}
  \caption{Efficiency by particle species for the two trackers. The error bars show the bootstrap uncertainty.}
  \label{fig:eff_species}
\end{figure}

\begin{table}[tbp]
  \caption{Efficiency by particle species, with the bootstrap uncertainty in percentage points.}
  \label{tab:species}
  \centering
  \begin{tabular}{lrcc}
    \hline\hline
    Species & Targets & Generic-trained (\%) & Signal-trained (\%) \\
    \hline
    muons & 2,289 & $97.52 \pm 0.31$ & $95.61 \pm 0.27$ \\
    electrons & 1,754 & $97.52 \pm 0.31$ & $96.16 \pm 0.32$ \\
    charged pions & 4,843 & $96.23 \pm 0.27$ & $94.76 \pm 0.27$ \\
    charged kaons & 1,533 & $96.36 \pm 0.37$ & $94.06 \pm 0.26$ \\
    protons & 889 & $94.04 \pm 0.29$ & $92.00 \pm 0.44$ \\
    \hline\hline
  \end{tabular}
\end{table}

Figure~\ref{fig:pileup} shows how the efficiency changes when pile-up is added to the test events.
At an average of 50 pile-up interactions, the generic-trained tracker finds 94.73\% of the targets
and the baseline 91.82\%. With 140 pile-up interactions, the efficiencies fall to 94.17\% and
90.69\%. At 200 interactions per crossing, the efficiencies are 93.72\% and 89.59\%, and the
fake-track rate of the generic-trained tracker rises to 1.12\%. Neither tracker was trained with
pile-up, so both see it here for the first time\footnote{The pile-up events are drawn from the same
generic sample as the training events but from disjoint event numbers.}. The generic-trained tracker
degrades more slowly, which suggests that it has learned the geometry of tracks rather than the
topology of the signal.

\begin{figure}[tbp]
  \centering
  \fbox{\parbox[c][0.45\linewidth][c]{0.8\linewidth}{\centering placeholder artwork: pileup}}
  \caption{Efficiency versus the average number of pile-up interactions. Neither tracker was trained with pile-up.}
  \label{fig:pileup}
\end{figure}

\subsection{Kinematics, signal points and pairs}
\label{sec:kinematics}

The efficiency also depends on the transverse momentum of the particle. Between 1 and 2~GeV, the
generic-trained tracker finds 93.47\% of the targets and the signal-trained one 90.66\%. From 2 to
5~GeV, the efficiencies are 94.28\% and 91.46\%. In the range from 5 to 10~GeV, they are 94.90\% and
92.62\%. Particles with transverse momenta from 10 to 20~GeV are found with efficiencies of 95.13\%
and 92.42\%. For the hardest particles, from 20 to 50~GeV, the efficiencies reach 95.77\% and
93.22\%. Soft particles curl more in the magnetic field, which makes their tracks harder to follow
from layer to layer. The lead of the generic-trained tracker is largest for soft particles, which
are common in generic events and rare in the signal.

The dependence on pseudorapidity is weak for both trackers. Both lose a little efficiency in the
transition between the barrel and the endcaps, where particles cross layers at shallow angles. The
difference between the trackers is the same in the barrel and in the endcaps within its uncertainty.
We therefore do not show it separately.

The headline comparison uses one signal point, but the conclusion should not depend on it. We
therefore repeated the comparison for four combinations of mass and lifetime of the heavy neutral
lepton. For a mass of 10~GeV and a decay length of 100~mm, the efficiencies are 95.41\% and 93.59\%.
At 15~GeV and 100~mm, which is the point the signal-trained tracker was trained on, they are 94.98\%
and 93.17\%. Doubling the decay length to 200~mm at 15~GeV gives 94.64\% and 92.34\%. The heaviest
and longest-lived point, 20~GeV and 300~mm, gives 94.16\% and 91.62\%. The signal-trained tracker
comes closest on its own training point, as expected. Away from it, its efficiency falls faster than
that of the generic-trained tracker.

A search would compare a learned tracker with the standard reconstruction, not only with another
learned tracker. The standard reconstruction with large-radius tracking finds 80.42\% of the
displaced targets at a fake-track rate of 3.36\%. It needs 435~ms per event for the large-radius
pass alone. Both learned trackers find many more displaced targets with far fewer fake tracks. This
comparison is less controlled than the main one, because the standard reconstruction differs in
every stage. We show it only to place the learned trackers on a familiar scale.

Many searches need both decay products of the long-lived particle, because a vertex fit combines
them~\cite{Fischer2022Vertexing}. The generic-trained tracker finds both daughters in 90.27\% of the
signal events, against 84.29\% for the signal-trained tracker. The pair efficiency is roughly the
square of the single-particle efficiency, so small gains per particle become larger gains per event.
The daughters of one decay are often close to each other, and a tracker can merge them into one
track. Merged pairs are rare for both trackers and slightly rarer for the generic-trained one.

\subsection{Failure modes}
\label{sec:failures}

We inspected the displaced targets that one tracker finds and the other misses. Most targets missed
only by the signal-trained tracker have a first hit in the strip layers. In these cases the
signal-trained network assigns low scores to the edges between the first two hits. The
generic-trained network gives the same edges high scores, because it has seen many tracks that start
in the strips from material interactions. Targets missed only by the generic-trained tracker are
fewer and show no common pattern.

Some failures are shared by both trackers. Particles that cross a gap between modules lose a hit and
may be split into two short tracks. Particles that scatter strongly in the material change direction
and are followed only up to the scattering point. Both failure modes depend on the detector and the
graph, not on the training, so they cancel in the comparison. They set the ceiling that neither
tracker can pass without a better graph.

The efficiency averages over particles, but a search selects events. We therefore also counted,
event by event, how many displaced targets each tracker finds. In most signal events both trackers
find the same displaced particles. Where they differ, the generic-trained tracker usually finds one
more particle than the signal-trained tracker. Events in which the signal-trained tracker finds more
are less common and contain mostly soft particles near the edge of the acceptance.

The training size matters for both trackers. With half of the training events, both trackers lose
some efficiency, and the difference between them stays about the same. With a quarter of the
training events, the signal-trained tracker loses more, and the difference grows. The
generic-trained tracker is thus less sensitive to the amount of training data. This is consistent
with the larger number of particles per generic event.

Prompt particles provide a control for the whole comparison. On prompt targets in the signal events,
the two trackers agree to within their uncertainties. The same holds for prompt targets in held-out
generic events. The difference on displaced targets is therefore not a general offset between the
trackers. It appears only for the particles whose production point lies away from the beam line.
This pattern is what one expects if the signal-trained tracker has learned the topology of one decay
rather than the geometry of tracks.

\subsection{Resolution, graph and timing}
\label{sec:resolution}

Figure~\ref{fig:resolution} compares the track parameters with the true values. The transverse
impact parameter is measured with a resolution of 0.018~mm, and the longitudinal one with 0.069~mm.
The relative transverse-momentum resolution is 1.93\%. The azimuthal angle is measured to 0.41~mrad.
These resolutions are the same for both trackers within their uncertainties, because they come from
the hits, not from the network.

\begin{figure}[tbp]
  \centering
  \fbox{\parbox[c][0.45\linewidth][c]{0.8\linewidth}{\centering placeholder artwork: resolution}}
  \caption{Resolution of the track parameters for matched displaced tracks. The two trackers overlap within their uncertainties.}
  \label{fig:resolution}
\end{figure}

Figure~\ref{fig:graph} shows the edge efficiency of the graph as a function of the layer pair. A
perfect edge classifier on the same graph would reach an efficiency of 99.35\%; this is the ceiling
of the pipeline. The generic-trained tracker closes most of the distance between the baseline and
this ceiling. The remaining loss comes from edges that the module map never proposes. Improving the
graph would raise both trackers by the same amount, so it would not change the comparison.

\begin{figure}[tbp]
  \centering
  \fbox{\parbox[c][0.45\linewidth][c]{0.8\linewidth}{\centering placeholder artwork: graph}}
  \caption{Edge efficiency of the constructed graph per pair of adjacent layers. The graph is the same for both trackers.}
  \label{fig:graph}
\end{figure}

Figure~\ref{fig:timing} shows the time spent in each stage of the pipeline on one graphics
processor. Graph construction takes 56~ms per event, the edge classifier 179~ms and track building
24~ms\footnote{The timing excludes reading the events from disk and writing the tracks.}. The total
of 259~ms compares with 296~ms for the standard reconstruction with large-radius tracking enabled.
The two trackers take the same time, because they share the architecture. The timing would allow the
pipeline to run in a software trigger, but we have not optimized it for that purpose.

\begin{figure}[tbp]
  \centering
  \fbox{\parbox[c][0.45\linewidth][c]{0.8\linewidth}{\centering placeholder artwork: timing}}
  \caption{Time per event spent in each stage of the pipeline, on one graphics processor. The standard reconstruction is shown for comparison.}
  \label{fig:timing}
\end{figure}

Figure~\ref{fig:ablation} shows which parts of the network carry the result. Removing the edge
features lowers the efficiency of the generic-trained tracker to 94.78\%. With four message-passing
steps instead of eight, the efficiency drops to 93.96\%. Without the curvature term in the loss, the
efficiency is 94.94\%. The number of message-passing steps matters most, because a displaced track
needs more steps to connect its first hit to the rest. None of the variants reverses the comparison
with the signal-trained tracker.

\begin{figure}[tbp]
  \centering
  \fbox{\parbox[c][0.45\linewidth][c]{0.8\linewidth}{\centering placeholder artwork: ablation}}
  \caption{Efficiency of the generic-trained tracker for variants of the network. The dashed line shows the signal-trained tracker.}
  \label{fig:ablation}
\end{figure}

\subsection{Checks and interpretation}
\label{sec:checks}

We checked that the result does not depend on details of the analysis. Changing the edge threshold
within the range allowed by the validation events moves both efficiencies by the same amount. Using
single instead of double majority raises both efficiencies and leaves the difference unchanged
within its uncertainty. Evaluating on the training process of each tracker shows that both reach a
similar efficiency on their own kind of events. The advantage of the generic-trained tracker
therefore appears only on events that neither tracker was trained on. Each check is repeated for
every training seed, and Appendix~\ref{app:seeds} shows the spread.

Why does the tracker that never saw the signal do better on it? The top-quark pair events contain
many more particles per event than the signal events. The generic-trained network therefore sees
many more examples of tracks, including the few displaced tracks from heavy-flavour decays. The
signal-trained network sees fewer tracks, and most of them share the kinematics of one decay. It
learns that decay well and generalizes less to the different lifetime and mass of the test sample.
This interpretation is consistent with the pile-up study in Fig.~\ref{fig:pileup}, where the
generic-trained tracker degrades more slowly.

Our study has three limitations. The detector is a generic model, and a real detector has dead
modules, misalignment and a more complex material map. The signal is one family of models, and other
long-lived signatures may behave differently. The test events have no pile-up in the headline
comparison, and the pile-up study uses a simplified overlay. None of these limitations favours one
of the trackers, so we expect the comparison to survive them. The numbers in
Table~\ref{tab:headline} and Table~\ref{tab:species} are therefore statements about this simulation
only.

\section{Discussion}
\label{sec:discussion}

The result has a practical consequence for experiments. Signal samples for every model in a search
programme are expensive to simulate, and their number grows with every new idea. Generic samples are
simulated anyway, in large numbers, for every physics analysis. If a tracker trained on generic
samples is at least as good on displaced signals, the expensive signal samples are needed only for
the evaluation. This makes learned tracking easier to adopt and easier to validate.

It also simplifies the statistical treatment of a search. When the tracker is trained on the signal,
the efficiency for a new model must be studied with care, because the tracker may favour the trained
model. A generic-trained tracker has no such preference, so its efficiency for a new model can be
estimated with the same procedure as for any other selection. The uncertainty of that estimate then
comes from the simulation of the detector, not from the training.

We do not claim that generic training is always better. A tracker trained on a large and varied
signal sample might match or surpass the generic-trained one. Our signal sample was as large as the
generic one in events but much smaller in particles. A fair comparison at equal numbers of training
particles is a natural extension of this study. The subset study in the appendix gives a first
indication that the generic events would still win.

Several open questions remain. We have not studied how the result changes with the geometry of the
detector. A detector with fewer layers gives a learned tracker less information per particle and may
change the balance between the trackers. We have also not studied tracks from particles that are not
charged leptons or hadrons, such as long-lived charged particles that decay inside the tracker. Such
particles leave short tracks that end abruptly, which neither training sample contains in large
numbers. Each of these questions can be answered with the same controlled comparison.


\section{Conclusions}
\label{sec:concl}
We have shown that a graph tracker trained without any signal finds more displaced particles than a classical tracker with a dedicated large-radius pass.
The gain is 13.54~$\pm$~0.31 points over all signal samples, and it is largest far from the beam line (Sec.~\ref{sec:radius}).
The fake-track rate is lower as well, and the time per event is shorter (Sec.~\ref{sec:timing}).
The principal limitation is that all results come from one simulated detector, mostly without pileup.
The next step is to repeat the test with the full pileup of the high-luminosity programme.

\begin{acknowledgments}
We thank the authors of the open detector simulation for their tools.
\end{acknowledgments}

% Appendix (generated by scripts/gen-large-fixture.mjs; synthetic text).
\appendix
\section{Training seeds}
\label{app:seeds}
A simpler filter is easier to deploy in the trigger. Every number in this paper is the mean over the 3 training seeds. The three training seeds give the same ranking in every bin. The learning rate were fixed on the validation split. The independent training seeds give the same ranking in every bin. The number of message-passing steps were fixed on the validation split. One simpler network needs more memory per event.

\begin{figure}[tbp]
  \centering
  \fbox{\parbox[c][0.28\textheight][c]{0.8\linewidth}{\centering placeholder artwork (fig:seed-spread)}}
  \caption{Efficiency versus production radius for each training seed. The spread between the seeds enters the uncertainty of Eq.~(\ref{eq:eff}).}
  \label{fig:seed-spread}
\end{figure}

The three training seeds change the efficiency by a small amount. Figure~\ref{fig:seed-spread} shows
each seed separately. A slower network gives a lower efficiency at large radius. Varying the
material in the simulation moves both trackers in the same direction. Varying the alignment of the
pixel layers changes the gain less than its statistical uncertainty. The independent training seeds
differ mostly at large radius. The three training seeds enter the uncertainty as a spread. Varying
the alignment of the pixel layers moves both trackers in the same direction.

The independent training seeds enter the uncertainty as a spread.
One smaller graph is easier to deploy in the trigger.
Scaling the pileup overlay affects the baseline more than the graph tracker.
The hyperparameters were not tuned on any signal sample.
Varying the pileup overlay changes the gain less than its statistical uncertainty.
A smaller network needs more memory per event.
The hyperparameters were fixed on the validation split.
A smaller network gives a lower efficiency at large radius.

\section{Material and alignment}
\label{app:material}
Removing the alignment of the pixel layers changes the gain less than its statistical uncertainty. The pileup overlay is the largest effect, at 0.25 points. Scaling the alignment of the pixel layers leaves the ranking unchanged. One smaller network gives a lower efficiency at large radius. The number of message-passing steps are the same for every seed. One smaller filter gives a lower efficiency at large radius. The learning rate were not tuned on any signal sample. A simpler network gives a lower efficiency at large radius.

\begin{figure}[tbp]
  \centering
  \fbox{\parbox[c][0.28\textheight][c]{0.8\linewidth}{\centering placeholder artwork (fig:material)}}
  \caption{Material in the tracker in units of radiation length, versus pseudorapidity. The variation used for the systematic effect scales every layer by the same factor.}
  \label{fig:material}
\end{figure}

One slower filter is easier to deploy in the trigger. Figure~\ref{fig:material} shows the material
budget. A smaller graph gives a lower efficiency at large radius. A slower network is easier to
deploy in the trigger. Varying the material in the simulation leaves the ranking unchanged. One
slower graph gives a lower efficiency at large radius. Removing the pileup overlay moves both
trackers in the same direction. The hyperparameters are the same for every seed.

\section{Hyperparameters}
\label{app:hyper}
The independent training seeds change the efficiency by a small amount.
The network has a hidden dimension of 128 and runs 8 message-passing steps.
A smaller filter needs more memory per event.
Training runs for 60 epochs with a learning rate of 0.0005.
Scaling the pileup overlay moves both trackers in the same direction.
One simpler graph gives a lower efficiency at large radius.
A smaller graph needs more memory per event.

\begin{table}[tbp]
  \caption{Hyperparameters of the edge classifier, fixed on the validation split.}
  \label{tab:hyper}
  \begin{tabular}{lc}
    \hline
    Parameter & Value \\
    Hidden dimension & 128 \\
    Message-passing steps & 8 \\
    Learning rate & 0.0005 \\
    Events per batch & 4 \\
    Epochs & 60 \\
    \hline
  \end{tabular}
\end{table}

Varying the material in the simulation affects the baseline more than the graph tracker. Scaling the material in the simulation leaves the ranking unchanged. Varying the alignment of the pixel layers leaves the ranking unchanged. A slower graph needs more memory per event. The number of message-passing steps were not tuned on any signal sample. A simpler graph needs more memory per event. The hidden dimension are the same for every seed. One smaller filter needs more memory per event.


\bibliography{refs}

\end{document}
